% !TEX program = xelatex
% Lernplatt - by Can Siebert
% Kurs 71 (Dig 42) - Tag 10 - Aufgabenheft
% Skalierbare digitale Prozessarchitektur: Microservices, Kubernetes, TDD/DevOps, Operating Model
%
% Self-contained xelatex source. Kein biblatex, keine externen Bilder.

\documentclass[11pt,a4paper,oneside]{article}

% --- Encoding & Sprache --------------------------------------------------
\usepackage{polyglossia}
\setdefaultlanguage{german}

% --- Geometrie -----------------------------------------------------------
\usepackage[a4paper,margin=2.5cm]{geometry}

% --- Fonts ---------------------------------------------------------------
\usepackage{fontspec}
\defaultfontfeatures{Ligatures=TeX}

\IfFontExistsTF{Charter}{%
  \setmainfont{Charter}%
}{%
  \IfFontExistsTF{Bitstream Charter}{%
    \setmainfont{Bitstream Charter}%
  }{%
    \setmainfont{TeX Gyre Schola}%
  }%
}

\IfFontExistsTF{Source Sans 3}{%
  \setsansfont{Source Sans 3}%
}{%
  \IfFontExistsTF{Source Sans Pro}{%
    \setsansfont{Source Sans Pro}%
  }{%
    \setsansfont{TeX Gyre Heros}%
  }%
}

\newcommand{\caveatfontfallback}{\itshape}
\IfFontExistsTF{Caveat}{%
  \newfontfamily\caveatfont{Caveat}[Scale=1.15]%
}{%
  \newcommand\caveatfont{\caveatfontfallback}%
}

\setmonofont{Latin Modern Mono}

% --- Mikro-Typografie ----------------------------------------------------
\usepackage{microtype}
\linespread{1.15}

% --- Farben & Boxen ------------------------------------------------------
\usepackage{xcolor}
\definecolor{lpInk}{RGB}{20,28,38}
% Kurs-71 Akzent: orange (#9d4a1a)
\definecolor{lpAccent}{RGB}{157,74,26}
\definecolor{lpAccentLight}{RGB}{246,228,212}
\definecolor{lpAmber}{RGB}{222,138,30}
\definecolor{lpAmberLight}{RGB}{255,243,217}
\definecolor{lpAmberBorder}{RGB}{196,118,18}
\definecolor{lpGray}{RGB}{96,104,114}
\definecolor{lpGrayLight}{RGB}{238,240,243}
\definecolor{lpGreen}{RGB}{34,120,72}
\definecolor{lpGreenLight}{RGB}{222,238,228}
\definecolor{lpGreenBorder}{RGB}{24,96,58}

% Niveau-Farben
\definecolor{lpL1}{RGB}{34,120,72}      % L1 = gruen (reproduktion)
\definecolor{lpL2}{RGB}{22,90,140}      % L2 = blau (anwendung)
\definecolor{lpL3}{RGB}{170,42,52}      % L3 = rot (management)

\usepackage[most]{tcolorbox}
\tcbuselibrary{breakable,skins}

% Info-Box (orange-weiss) fuer Hinweise/Aufgabenstellung-Rahmen
\newtcolorbox{infobox}[1][Hinweis]{%
  enhanced, breakable,
  colback=lpAccentLight,
  colframe=lpAccent,
  boxrule=0.6pt,
  arc=2pt,
  left=10pt,right=10pt,top=8pt,bottom=8pt,
  fonttitle=\sffamily\bfseries\color{white},
  coltitle=white,
  title={#1},
  attach boxed title to top left={xshift=8pt,yshift=-8pt},
  boxed title style={size=small, colback=lpAccent, colframe=lpAccent, boxrule=0.4pt, arc=1pt},
  top=14pt
}

% Antwort-Bereich: leerer Kasten mit Punktlinien
\newtcolorbox{antwortbox}[1][Ihre Antwort]{%
  enhanced, breakable,
  colback=white,
  colframe=lpGray,
  boxrule=0.4pt,
  arc=2pt,
  left=10pt,right=10pt,top=8pt,bottom=8pt,
  fonttitle=\sffamily\bfseries\color{white},
  coltitle=white,
  title={#1},
  attach boxed title to top left={xshift=8pt,yshift=-8pt},
  boxed title style={size=small, colback=lpGray, colframe=lpGray, boxrule=0.4pt, arc=1pt},
  top=14pt
}

% --- Listen --------------------------------------------------------------
\usepackage{enumitem}
\setlist{itemsep=2pt, topsep=4pt, parsep=0pt}
\setlist[itemize,1]{label={\textbullet},leftmargin=1.6em}
\setlist[itemize,2]{label={\textendash},leftmargin=1.4em}
\setlist[enumerate,1]{label=\arabic*., leftmargin=1.8em}

% --- Tabellen ------------------------------------------------------------
\usepackage{tabularx}
\usepackage{array}
\usepackage{booktabs}
\usepackage{longtable}

% --- Kopf-/Fusszeilen ----------------------------------------------------
\usepackage{fancyhdr}
\usepackage{lastpage}
\pagestyle{fancy}
\fancyhf{}
\renewcommand{\headrulewidth}{0.3pt}
\renewcommand{\footrulewidth}{0pt}
\fancyhead[L]{\footnotesize\sffamily\color{lpGray}Aufgabenheft \textperiodcentered{} Kurs 71 \textperiodcentered{} Tag 10}
\fancyhead[R]{\footnotesize\sffamily\color{lpGray}Microservices, Kubernetes \& DevOps}
\fancyfoot[C]{\footnotesize\sffamily\color{lpGray}\thepage{} / \pageref{LastPage}}

\fancypagestyle{plain}{%
  \fancyhf{}%
  \renewcommand{\headrulewidth}{0pt}%
  \fancyfoot[C]{\footnotesize\sffamily\color{lpGray}\thepage{} / \pageref{LastPage}}%
}

% --- Section-Styling -----------------------------------------------------
\usepackage[explicit]{titlesec}

\titleformat{\section}[block]
  {\sffamily\Large\bfseries\color{lpAccent}}
  {\thesection.}{0.6em}{#1}
  [{\vspace{2pt}\color{lpAccent}\titlerule[0.6pt]}]

\titleformat{\subsection}[block]
  {\sffamily\large\bfseries\color{lpInk}}
  {\thesubsection}{0.6em}{#1}

\titlespacing*{\section}{0pt}{14pt}{8pt}
\titlespacing*{\subsection}{0pt}{10pt}{4pt}

% --- Hyperlinks ----------------------------------------------------------
\usepackage[hidelinks,unicode]{hyperref}

% --- Helfer --------------------------------------------------------------
\newcommand{\akzent}[1]{{\caveatfont\color{lpAmber}#1}}
\newcommand{\fachbegriff}[1]{\textbf{#1}}

% Niveau-Badge
\newcommand{\niveaubadge}[2]{%
  \tcbox[on line, colback=#1, colframe=#1, coltext=white,
         boxrule=0pt, arc=2pt, left=5pt,right=5pt,top=1pt,bottom=1pt,
         nobeforeafter]{\sffamily\bfseries\footnotesize #2}%
}
\newcommand{\nivLi}{\niveaubadge{lpL1}{L1\,\textbullet\,Wissen}}
\newcommand{\nivLii}{\niveaubadge{lpL2}{L2\,\textbullet\,Anwendung}}
\newcommand{\nivLiii}{\niveaubadge{lpL3}{L3\,\textbullet\,Management}}

\newcommand{\zeit}[1]{%
  \tcbox[on line, colback=lpGrayLight, colframe=lpGrayLight, coltext=lpInk,
         boxrule=0pt, arc=2pt, left=5pt,right=5pt,top=1pt,bottom=1pt,
         nobeforeafter]{\sffamily\footnotesize\textbf{$\sim$\,#1}}%
}

\newcommand{\aufgabenkopf}[4]{%
  \par\vspace{6pt}
  \noindent{\sffamily\Large\bfseries\color{lpAccent}Aufgabe #1 \textendash{} #2}\par
  \vspace{2pt}
  \noindent{\color{lpAccent}\rule{\linewidth}{0.6pt}}\par
  \vspace{4pt}
  \noindent #3 \hfill \zeit{#4}\par
  \vspace{6pt}
}

\newcommand{\ytquery}[1]{%
  \par\vspace{4pt}
  \noindent{\footnotesize\sffamily\color{lpGray}\textbf{YouTube-Suchquery:} \texttt{#1}}\par
}

\newcommand{\antwortlinien}[1]{%
  \par\vspace{6pt}
  \foreach \x in {1,...,#1}{%
    \noindent\makebox[\linewidth]{\dotfill}\\[6pt]
  }
}
\usepackage{pgffor}

% =========================================================================
\begin{document}

% --- COVER ---------------------------------------------------------------
\begin{titlepage}
  \pagestyle{empty}
  \vspace*{1.5cm}
  \noindent{\sffamily\footnotesize\color{lpGray}LERNPLATT \textperiodcentered{} BY CAN SIEBERT}\\[2pt]
  \noindent{\color{lpAccent}\rule{\linewidth}{1.2pt}}\\[4pt]
  \noindent{\sffamily\footnotesize\color{lpGray}AUFGABENHEFT \textperiodcentered{} MODUL 4 \textperiodcentered{} TAG 10 VON 16 \textperiodcentered{} KURS 71 (Dig 42) \textperiodcentered{} 11.\,Juni 2026}

  \vspace{3cm}
  {\sffamily\fontsize{30}{34}\selectfont\bfseries\color{lpInk}Aufgabenheft\\Tag 10\par}

  \vspace{0.7cm}
  {\sffamily\large\color{lpGray}Skalierbare digitale Prozess\-architektur \textendash{} Microservices,\\Kubernetes, TDD/DevOps und Operating\\Model in elf Aufgaben.\par}

  \vspace{1.5cm}
  {\caveatfont\LARGE\color{lpAmber}11 Aufgaben \textperiodcentered{} L1 \textperiodcentered{} L2 \textperiodcentered{} L3}\par

  \vfill

  \noindent\begin{minipage}{0.55\linewidth}
    {\sffamily\footnotesize\color{lpGray}Niveau-Mix:}\\
    {\sffamily\small\color{lpInk}3\,\textsf{L1} \textperiodcentered{} 5\,\textsf{L2} \textperiodcentered{} 3\,\textsf{L3}}\\[10pt]
    {\sffamily\footnotesize\color{lpGray}Bearbeitung:}\\
    {\sffamily\small\color{lpInk}Selbstbearbeitung im Lernpfad}
  \end{minipage}\hfill
  \begin{minipage}{0.4\linewidth}
    \raggedleft
    {\sffamily\footnotesize\color{lpGray}Dozent:}\\
    {\sffamily\small\color{lpInk}Can Siebert}\\[10pt]
    {\sffamily\footnotesize\color{lpGray}Begleitend:}\\
    {\sffamily\small\color{lpInk}Lehrskript \& L\"osungsheft}
  \end{minipage}

  \vspace{0.5cm}
  \noindent{\color{lpAccent}\rule{\linewidth}{0.6pt}}
\end{titlepage}

% --- ANLEITUNG -----------------------------------------------------------
\section*{So arbeiten Sie mit diesem Heft}
\thispagestyle{plain}

\noindent Dieses Aufgabenheft enth\"alt elf Aufgaben zu Tag 10 \textendash{} \emph{Skalierbare digitale Prozessarchitektur}: Microservices, Container/Kubernetes, TDD und Teststrategien sowie DevOps-Strategien inklusive Operating Model. Die Aufgaben sind in drei Niveaus gegliedert:

\begin{itemize}
  \item \nivLi{} \textendash{} \fachbegriff{Wissen.} Begriffe, Servicegrenzen, Pipeline-Logik, Testpyramide. 5\,--\,10 Minuten pro Aufgabe.
  \item \nivLii{} \textendash{} \fachbegriff{Anwendung.} Servicegrenzen ziehen, Betriebskonzept entwerfen, SLOs definieren, RACI bauen. 15\,--\,30 Minuten.
  \item \nivLiii{} \textendash{} \fachbegriff{Management.} Operating Model, Entscheidungsarchitektur, Reliability-Framework. 30\,--\,45 Minuten.
\end{itemize}

\noindent Jede Aufgabe enth\"alt:

\begin{itemize}
  \item eine Niveau-Markierung und einen Zeit-Richtwert,
  \item den Aufgabentext (h\"aufig mit Fallbeschreibung),
  \item einen Antwort-Freiraum f\"ur Ihre L\"osung,
  \item eine \emph{YouTube-Suchquery} f\"ur den begleitenden Lernpfad.
\end{itemize}

\begin{infobox}[Hinweis]
Die Musterl\"osungen finden Sie im separaten \emph{L\"osungsheft Tag 10}. Versuchen Sie jede Aufgabe zuerst eigenst\"andig. Der Mehrwert entsteht im Entwerfen und Argumentieren \textendash{} nicht im Abschreiben.
\end{infobox}

\clearpage

% =========================================================================
% L1 AUFGABEN
% =========================================================================
\section*{Niveau L1 \textendash{} Wissen \& Reproduktion}
\addcontentsline{toc}{section}{L1 -- Wissen}

% --- AUFGABE 1 -----------------------------------------------------------
\aufgabenkopf{1}{Microservices \textendash{} Eignung, Risiken und Servicegrenze}{\nivLi}{8\,min}

\noindent Erl\"autern Sie in eigenen Worten, was \emph{Microservices} sind und worin sich ihr Vorteil gegen\"uber einem Monolithen begr\"undet. Beantworten Sie konkret:

\begin{enumerate}
  \item Nennen Sie \textbf{drei Nutzenversprechen} von Microservices (z.\,B.\ unabh\"angige Deployments, gezielte Skalierung, klare Verantwortungsgrenzen).
  \item Nennen Sie \textbf{drei Risiken}, die mit Microservices typischerweise einhergehen (z.\,B.\ Komplexit\"at, verteilte Fehler, Observability-Bedarf).
  \item Erl\"autern Sie in einem Satz, was eine \emph{gute Servicegrenze} ausmacht \textendash{} und warum sie \emph{nicht} entlang von Datenbanktabellen, sondern entlang von \emph{Gesch\"aftsf\"ahigkeiten (Capabilities)} verlaufen sollte.
\end{enumerate}

\ytquery{Microservices Vorteile Risiken Servicegrenze Business Capability erklaert}

\antwortlinien{14}

\clearpage

% --- AUFGABE 2 -----------------------------------------------------------
\aufgabenkopf{2}{Container \& Kubernetes \textendash{} Kernbegriffe}{\nivLi}{8\,min}

\noindent Ordnen Sie die folgenden Kernbegriffe der Container- und Kubernetes-Welt jeweils ihrer korrekten Bedeutung zu. Schreiben Sie zu jedem Begriff einen kurzen, eigenen Erkl\"arungssatz \textendash{} so, dass eine fachfremde Kollegin den Begriff versteht.

\medskip
\noindent\textbf{Begriffe:}

\begin{enumerate}
  \item \textbf{Container}
  \item \textbf{Orchestrierung}
  \item \textbf{Self-Healing}
  \item \textbf{Rollback}
  \item \textbf{Service-Discovery}
  \item \textbf{Blue/Green-Deployment}
  \item \textbf{Canary-Release}
  \item \textbf{Ressourcen\-limits (Requests/Limits)}
\end{enumerate}

\noindent Markieren Sie zus\"atzlich zwei Begriffe, deren \emph{Auslassen} im Betrieb am schnellsten zu Schaden f\"uhrt \textendash{} und begr\"unden Sie kurz, warum.

\ytquery{Kubernetes Grundlagen Container Orchestrierung Blue Green Canary einfach erklaert}

\antwortlinien{14}

\clearpage

% --- AUFGABE 3 -----------------------------------------------------------
\aufgabenkopf{3}{Testpyramide \& Quality Gates}{\nivLi}{10\,min}

\noindent Erl\"autern Sie das Prinzip der \emph{Testpyramide} und erstellen Sie eine \textbf{Vergleichstabelle} f\"ur die drei Teststufen \emph{Unit}, \emph{Integration} und \emph{End-to-End (E2E)}.

\medskip
\noindent\begin{tabularx}{\linewidth}{@{}p{3.6cm}|X|X|X@{}}
\toprule
\textbf{Eigenschaft} & \textbf{Unit-Test} & \textbf{Integrationstest} & \textbf{E2E-Test} \\
\midrule
Was wird getestet?  & & & \\[16pt]
\midrule
Geschwindigkeit     & & & \\[16pt]
\midrule
Kosten / Fragilit\"at & & & \\[16pt]
\midrule
Anteil in Pyramide  & & & \\[16pt]
\midrule
Typischer Fehlerfund & & & \\[16pt]
\bottomrule
\end{tabularx}

\smallskip
\noindent Erg\"anzen Sie unter der Tabelle in zwei S\"atzen: Was sind \emph{Quality Gates}, und welche drei Pr\"ufungen w\"urden Sie als nicht verhandelbares Gate vor einem Produktiv-Release verlangen?

\ytquery{Testpyramide TDD Red Green Refactor Quality Gates erklaert}

\antwortlinien{8}

\clearpage

% =========================================================================
% L2 AUFGABEN
% =========================================================================
\section*{Niveau L2 \textendash{} Anwendung \& Transfer}
\addcontentsline{toc}{section}{L2 -- Anwendung}

% --- AUFGABE 4 -----------------------------------------------------------
\aufgabenkopf{4}{Service-Schnitt aus Prozessbeschreibung \textendash{} Order-to-Cash}{\nivLii}{25\,min}

\begin{infobox}[Ausgangslage: Order-to-Cash]
\textbf{Prozess:} Bestellung \textrightarrow{} Verf\"ugbarkeit pr\"ufen \textrightarrow{} Preis/Angebot \textrightarrow{} Auftrag best\"atigen \textrightarrow{} Versand \textrightarrow{} Rechnung \textrightarrow{} Zahlung \textrightarrow{} Reklamation (ggf.) \textrightarrow{} Reporting.

\smallskip
\textbf{Ausgangszustand:} Heute ein Monolith, Releases alle 6 Wochen, viele Abh\"angigkeiten zwischen Modulen.

\smallskip
\textbf{Cheat Sheet:} Service = klarer Zweck + klarer Owner + Schnittstelle (Input/Output). Gute Grenze = m\"oglichst wenig ``Hin und Her'' zwischen Services (geringe Kopplung).
\end{infobox}

\noindent\textbf{Arbeitsauftrag:}

\begin{enumerate}
  \item Schneiden Sie den Prozess in \textbf{5\,--\,7 Service-Kandidaten} (z.\,B.\ \emph{Order}, \emph{Inventory}, \emph{Shipping}, \emph{Billing}, \emph{Payments}, \emph{Complaints}).
  \item F\"ur \textbf{drei} dieser Services definieren Sie einen \emph{Service-Steckbrief} mit:
    \begin{itemize}
      \item \emph{Zweck} (1 Satz),
      \item \emph{Input} (vier Felder) und \emph{Output} (vier Felder),
      \item \emph{Owner-Rolle},
      \item \emph{kritischster Fehlerfall} (1 Satz).
    \end{itemize}
  \item Entscheiden Sie unter Unsicherheit: Welche \emph{eine} Service-Grenze ist am riskantesten (weil Daten- oder Ownership-Verh\"altnisse unklar sind)? Definieren Sie einen \textbf{Guardrail}, der das Risiko begrenzt (z.\,B.\ Anti-Corruption-Layer, fachliche Eskalation, gemeinsamer Datenvertrag).
\end{enumerate}

\ytquery{Microservices Service Schnitt Order to Cash Domain Driven Design Capability}

\antwortlinien{22}

\clearpage

% --- AUFGABE 5 -----------------------------------------------------------
\aufgabenkopf{5}{Release-Flow \& Testpyramide als Sicherheitsnetz}{\nivLii}{25\,min}

\begin{infobox}[Mini-Szenario]
\textbf{Ziel:} W\"ochentliches Release. \textbf{Sorge:} Ausf\"alle in der Produktion. \textbf{Pflicht:} Der Kernprozess \emph{Rechnung / Zahlung} muss auditierbar bleiben.

\smallskip
\textbf{Cheat Sheet:} Deployment braucht Rollout-Plan, Rollback, Monitoring/Alerts und Verantwortliche. Tests gliedern sich in Unit (schnell), Integration (Schnittstellen), E2E (Gesch\"aftsfluss).
\end{infobox}

\noindent\textbf{Arbeitsauftrag:}

\begin{enumerate}
  \item Entwerfen Sie einen \textbf{minimalen Release-Flow} in maximal \emph{sieben Schritten} von ``\"Anderung'' bis ``Prod'' inklusive zweier \textbf{Quality Gates}.
  \item Skizzieren Sie eine \textbf{Testpyramide f\"ur drei Services} (z.\,B.\ Order, Billing, Payments). Geben Sie f\"ur jeden Service grob an, wie viele Tests welcher Art Sie ansetzen.
  \item Definieren Sie \textbf{zwei Betriebskennzahlen} (SLO-orientiert), z.\,B.\ \emph{Verf\"ugbarkeit}, \emph{Fehlerquote}, \emph{Latenz}, \emph{MTTR} \textendash{} jeweils mit \emph{Owner} und \emph{Zielwert}.
  \item Entscheidung unter Unsicherheit: Wo akzeptieren Sie im \textbf{MVP} bewusst Rest-Risiko (z.\,B.\ weniger E2E-Tests)? Wie kompensieren Sie es (Canary, Feature Toggle, manuelle Kontrolle)? Dokumentieren Sie das in einem kurzen \emph{Risikoplan}.
\end{enumerate}

\ytquery{Release Flow Quality Gates Canary Feature Toggle SLO Beispiel}

\antwortlinien{20}

\clearpage

% --- AUFGABE 6 -----------------------------------------------------------
\aufgabenkopf{6}{Fallstudie ``InsureFast'' \textendash{} Servicegrenzen \& Betriebskonzept}{\nivLii}{30\,min}

\begin{infobox}[Fallstudie: InsureFast]
\textbf{Unternehmen:} ``InsureFast'' \textendash{} Versicherer mit digitalem Abschluss und Schadenmeldung.

\smallskip
\textbf{Problem:} Release-Zyklus 6\,--\,8 Wochen; h\"aufige Produktionsprobleme; der Kundenprozess \emph{Schaden melden} bricht bei Lastspitzen; Audit fordert Nachvollziehbarkeit (\emph{wer hat wann was entschieden?}).

\smallskip
\textbf{Restriktionen / Unsicherheit:} 12 Wochen bis Saison-Peak, Budget 220\,k\,\euro{}, Legacy-DB, Security-Anforderungen, unklare Lastprofile, Zielkonflikt \emph{Speed vs.\ Stabilit\"at}.
\end{infobox}

\noindent\textbf{Arbeitsauftrag (Teil 1 von 2):}

\begin{enumerate}
  \item \textbf{Servicegrenzen (5\,--\,8):} Definieren Sie Servicegrenzen entlang der Prozesskette \emph{Schaden melden \textrightarrow{} pr\"ufen \textrightarrow{} entscheiden \textrightarrow{} auszahlen} (z.\,B.\ \emph{Claim Intake}, \emph{Policy/Customer}, \emph{Fraud Check}, \emph{Claim Assessment}, \emph{Payout}, \emph{Notifications}, \emph{Document Management}). Begr\"unden Sie eine Grenze, die nicht trivial ist.
  \item \textbf{Kubernetes-Betriebskonzept (konzeptionell):} Beschreiben Sie f\"unf Bausteine:
    \begin{itemize}
      \item Rollout-/Rollback-Strategie f\"ur kritische Services,
      \item Skalierungsregel (Wenn-Dann),
      \item Umgang mit Konfiguration und Secrets,
      \item Monitoring/Alerting,
      \item Ressourcenlimits.
    \end{itemize}
\end{enumerate}

\ytquery{Microservices Versicherung Schadenmeldung Kubernetes Skalierung Beispiel}

\antwortlinien{22}

\clearpage

% --- AUFGABE 7 -----------------------------------------------------------
\aufgabenkopf{7}{Fallstudie ``InsureFast'' \textendash{} Teststrategie \& Operating Model}{\nivLii}{30\,min}

\begin{infobox}[Fortsetzung: InsureFast]
\textbf{Bezug:} Diese Aufgabe baut auf Aufgabe 6 auf. Sie setzen die Fallstudie ``InsureFast'' fort und erg\"anzen die fehlenden Bausteine.
\end{infobox}

\noindent\textbf{Arbeitsauftrag (Teil 2 von 2):}

\begin{enumerate}
  \item \textbf{Teststrategie (TDD-Ansatz):} Definieren Sie f\"ur InsureFast eine Teststrategie mit:
    \begin{itemize}
      \item Unit-Tests f\"ur Kernlogik (welche Logik prim\"ar?),
      \item Integrationstests f\"ur Service-APIs (welche Schnittstellen?),
      \item E2E-Tests nur f\"ur \emph{zwei} Kernpfade (welche Pfade w\"ahlen Sie?).
    \end{itemize}
    Erg\"anzen Sie \emph{zwei Quality Gates}, die jeder Build durchlaufen muss \textendash{} und eine \emph{minimale} E2E-Absicherung.
  \item \textbf{Operating Model:} Skizzieren Sie:
    \begin{itemize}
      \item den Teamzuschnitt (Produktteams + Plattformteam + Security/Compliance),
      \item eine RACI-Aufteilung f\"ur Schnittstellen\"anderungen,
      \item das Incident-Handling (wer wird wann gerufen?),
      \item den Release-Freigabepfad.
    \end{itemize}
  \item \textbf{Zwei Entscheidungen unter Unsicherheit:}
    \begin{itemize}
      \item \emph{Datenstrategie:} Big-Bang-Migration vs.\ DB-bleibt + Anti-Corruption-Layer \textendash{} was w\"ahlen Sie und warum?
      \item \emph{Release-Frequenz:} W\"ochentlich f\"ur alle Services vs.\ gestuft (Kern w\"ochentlich/zweiw\"ochentlich) \textendash{} welche Variante?
    \end{itemize}
    Definieren Sie zu jeder Entscheidung ein \textbf{Fr\"uhwarnsignal}.
\end{enumerate}

\ytquery{Microservices Operating Model RACI Plattform Team Incident Management}

\antwortlinien{22}

\clearpage

% --- AUFGABE 8 -----------------------------------------------------------
\aufgabenkopf{8}{Anti-Muster ``Microservices-Spaghetti'' erkennen und korrigieren}{\nivLii}{25\,min}

\begin{infobox}[Vier Beispiele aus dem Alltag]
\textbf{A \textendash{} ``Service-Wildwuchs'':} 47 Services f\"ur einen Mittelst\"andler mit 60 Entwicklern. Jeder Service hat einen eigenen Datenstand, dieselben Kundendaten existieren mehrfach.\\[2pt]

\textbf{B \textendash{} ``Verteilter Monolith'':} F\"unf Services, aber jede \"Anderung erfordert \emph{koordinierte Releases} aller f\"unf \textendash{} weil sie sich gegenseitig synchron aufrufen.\\[2pt]

\textbf{C \textendash{} ``Observability als Nachgedanke'':} Nach jedem Incident dauert es im Schnitt zwei Stunden, bis \"uberhaupt klar ist, welcher Service den Fehler ausgel\"ost hat. Logs liegen verstreut in f\"unf Tools.\\[2pt]

\textbf{D \textendash{} ``Wir vertrauen unseren Tests'' ohne Rollback-\"Ubung:} Niemand hat in den letzten sechs Monaten einen Rollback in der Produktion gefahren. Im Incident-Fall wird improvisiert.
\end{infobox}

\noindent\textbf{Arbeitsauftrag:}

\noindent F\"ur jedes der vier Beispiele (A, B, C, D):

\begin{enumerate}
  \item Benennen Sie das prim\"are \textbf{Anti-Muster} (z.\,B.\ ``Verteilter Monolith'', ``Service-Wildwuchs'', ``Observability-Schuld'', ``Rollback-Atrophie'').
  \item Beschreiben Sie in einem Satz, \textbf{welches Prinzip} dabei verletzt wird (Servicegrenze entlang Capabilities, lose Kopplung, Betriebsdisziplin, you-build-it-you-run-it).
  \item Schlagen Sie eine \textbf{konkrete Korrektur} in 1\,--\,2 S\"atzen vor.
\end{enumerate}

\ytquery{Microservices Antipattern verteilter Monolith Observability erklaert}

\antwortlinien{18}

\clearpage

% =========================================================================
% L3 AUFGABEN
% =========================================================================
\section*{Niveau L3 \textendash{} Management \& Entscheidungsarchitektur}
\addcontentsline{toc}{section}{L3 -- Management}

% --- AUFGABE 9 -----------------------------------------------------------
\aufgabenkopf{9}{Reliability- \& Risk-Framework f\"ur eine wachsende Plattform}{\nivLiii}{35\,min}

\begin{infobox}[Ausgangslage: Plattform-Skalierung]
\textbf{Ihre Rolle:} Head of Platform einer Organisation mit acht Produktteams. Drei Teams releasen w\"ochentlich, f\"unf releasen alle 4\,--\,6 Wochen.

\smallskip
\textbf{Beobachtung:} Die Change Failure Rate liegt bei 18\,\% (Soll: $\leq$\,10\,\%). MTTR ist hoch (Median 2,5\,h). Es gibt keine einheitlichen SLOs.

\smallskip
\textbf{Zielkonflikt:} Geschwindigkeit der Produktteams vs.\ Stabilit\"at und Auditierbarkeit der Plattform.
\end{infobox}

\noindent\textbf{Arbeitsauftrag:}

\begin{enumerate}
  \item \textbf{SLO-Set entwickeln:} Definieren Sie ein Minimal-SLO-Set (drei SLOs) f\"ur kritische Services \textendash{} mit Schwellenwerten und Begr\"undung (Verf\"ugbarkeit, Latenz, Fehlerquote).
  \item \textbf{Error Budgets (konzeptionell):} Erkl\"aren Sie, wie ein Error-Budget-Mechanismus zwischen Speed und Reliability vermittelt. Definieren Sie eine Regel, was passiert, wenn das Budget aufgebraucht ist (Release-Stop? Plattform-Review?).
  \item \textbf{Risk-Framework:} Erstellen Sie eine Liste von f\"unf typischen Risiken (z.\,B.\ \emph{Servicegrenze falsch geschnitten}, \emph{Secret im Code}, \emph{Rollback nicht ge\"ubt}, \emph{Observability-L\"ucken}, \emph{Tooling-Lock-in}). Bewerten Sie jedes Risiko nach \emph{Impact} und \emph{Wahrscheinlichkeit} (jeweils 1\,--\,5) und benennen Sie pro Risiko \emph{eine} pr\"aventive Ma{\ss}nahme.
  \item \textbf{Anti-Gaming-Mechanik:} Wie verhindern Sie, dass Teams ihre Change Failure Rate sch\"onen (z.\,B.\ Fehler nicht als Change deklarieren, kleine Releases h\"aufen, Postmortems vermeiden)? Definieren Sie zwei Gegenkopplungen.
\end{enumerate}

\ytquery{SLO Error Budget Change Failure Rate MTTR SRE Framework Senior}

\antwortlinien{28}

\clearpage

% --- AUFGABE 10 ----------------------------------------------------------
\aufgabenkopf{10}{Governance leicht, aber wirksam}{\nivLiii}{30\,min}

\noindent\textbf{Diskussionsaufgabe.}

\noindent Microservices und DevOps verlocken zu maximaler Teamautonomie. Ohne Leitplanken entstehen schnell Schatten-Architekturen, inkonsistente Schnittstellen und Audit-Findings. Schwergewichtige Architektur-Boards hingegen werden umgangen oder verz\"ogern Releases. Hier ist die Management-Spannung.

\medskip
\noindent\textbf{Arbeitsauftrag:}

\begin{enumerate}
  \item \textbf{Architekturprinzipien (max.\ 8 Prinzipien):} Formulieren Sie acht Prinzipien, die in Ihrer Organisation als \emph{verbindlich} gelten sollen (z.\,B.\ \emph{Servicegrenzen folgen Capabilities}, \emph{kein Secret im Code}, \emph{jeder Service hat einen Owner}, \emph{API-Versionierung verpflichtend}). Schreiben Sie jedes Prinzip in einem Satz und ohne F\"ullw\"orter.
  \item \textbf{Schnittstellenstandards:} Definieren Sie zwei verbindliche Standards (z.\,B.\ \emph{OpenAPI-Spezifikation f\"ur jede REST-API}, \emph{Event-Schema-Registry f\"ur asynchrone Kommunikation}, \emph{Idempotenz bei schreibenden Operationen}).
  \item \textbf{Review-Regel:} Definieren Sie, wann ein \emph{Architektur-Review} verpflichtend ist \textendash{} und wann nicht. Faustregel: Review nur f\"ur \emph{teure} oder \emph{irreversible} Entscheidungen.
  \item \textbf{Eskalationspfad:} Wer entscheidet, wenn Produktteam und Plattformteam sich uneinig sind? Skizzieren Sie den Pfad in vier Stufen.
\end{enumerate}

\ytquery{Architecture Governance Principles Senior Lightweight Service Standard}

\antwortlinien{20}

\clearpage

% --- AUFGABE 11 ----------------------------------------------------------
\aufgabenkopf{11}{Transferprojekt \textendash{} Digitale Prozessarchitektur als Betriebs- und Entscheidungsarchitektur}{\nivLiii}{45\,min}

\begin{infobox}[Rolle und Auftrag]
\textbf{Ihre Rolle:} Chief Architect / Head of Platform / Chief Digital Officer.

\smallskip
\textbf{Auftrag:} In \textbf{16 Wochen} eine skalierbare digitale Prozessplattform etablieren, die w\"ochentlich releasen kann \emph{und} auditierbar bleibt (Microservices + Kubernetes + DevOps + Operating Model).

\smallskip
\textbf{Restriktionen / Unsicherheit:}
\begin{itemize}
  \item Legacy-Landschaft, Lastprofile unbekannt,
  \item Budget \textbf{350\,k\,\euro{}},
  \item mehrere Teams mit unterschiedlichem Reifegrad,
  \item Security- und Audit-Druck,
  \item Zeitdruck durch Peak-Saison.
\end{itemize}
\end{infobox}

\noindent\textbf{Arbeitsauftrag:} Entwickeln Sie eine \emph{Entscheidungsarchitektur} (keine Ma{\ss}nahmenliste) f\"ur die Plattform. Erarbeiten Sie schriftlich:

\begin{enumerate}
  \item \textbf{Top-10-Entscheidungsarchitektur:} Servicegrenzen, Datenhoheit, API-/Event-Standards, SLOs, Release-Policies, Security-Baselines, Observability-Standard, Incident-Governance, Ausnahmeprozess, Toolchain-Standardisierung. F\"ur jede Entscheidung: Wer entscheidet? Welche Optionen?
  \item \textbf{Operating Model \& RACI:} Produktteams vs.\ Plattform vs.\ Security/Compliance. \emph{Wer darf was \"andern?} Eskalationspfad bei Zielkonflikten (Speed vs.\ Risk).
  \item \textbf{Reliability- \& Risk-Framework:} SLOs, Error Budgets, Change Failure Rate, MTTR. Definieren Sie eine \emph{Anti-Gaming}-Kopplung.
  \item \textbf{Governance leicht, aber wirksam:} max.\ acht Architekturprinzipien, Review nur f\"ur teure/irreversible Entscheidungen, Standard-Templates (CI/CD, Helm/Deployment \textendash{} konzeptionell).
  \item \textbf{Roadmap:} MVP (8 Wochen) \textrightarrow{} Scale (8 Wochen) inklusive \emph{Messkonzept} (Outcomes, nicht Aktivit\"at).
  \item \textbf{Zwei Entscheidungen unter Unsicherheit (Pflicht):}
    \begin{itemize}
      \item \emph{Welche drei Services zuerst extrahieren?} Begr\"undung + verworfene Alternativen + Fr\"uhwarnsignale.
      \item \emph{Welche Release-Policy f\"ur Kernprozesse?} Trade-off Speed/Safety + Guardrails.
    \end{itemize}
\end{enumerate}

\noindent\textbf{Bewertungskriterien (Senior-Level):} Entscheidungsklarheit \textperiodcentered{} Anschlussf\"ahigkeit an Strategie \textperiodcentered{} Pragmatismus unter Restriktionen (16 Wochen, 350\,k\,\euro{}, Reifegrad-Heterogenit\"at) \textperiodcentered{} Risikoreife (sichtbare, nicht verdr\"angte Risiken).

\ytquery{Microservices Operating Model Decision Architecture Senior Platform CIO}

\antwortlinien{36}

\clearpage

% --- Abschluss -----------------------------------------------------------
\section*{Abschluss}
\thispagestyle{plain}

\noindent Sie haben alle elf Aufgaben bearbeitet. Vergleichen Sie nun Ihre Antworten mit dem \emph{L\"osungsheft Tag 10}.

\medskip
\begin{infobox}[Was Sie jetzt k\"onnen sollten]
\begin{itemize}
  \item Microservices nach Capability-Logik schneiden und Servicegrenzen begr\"unden.
  \item Container/Kubernetes-Konzepte (Orchestrierung, Self-Healing, Rollback, Service-Discovery) zuordnen.
  \item Eine Testpyramide entwerfen und Quality Gates als Sicherheitsnetz formulieren.
  \item Ein Kubernetes-Betriebskonzept entwerfen (Rollout, Skalierung, Konfig/Secrets, Monitoring).
  \item Ein Operating Model mit RACI f\"ur Produktteams, Plattform und Security entwickeln.
  \item Ein Reliability-Framework (SLOs, Error Budgets, MTTR) inklusive Anti-Gaming-Kopplung entwerfen.
  \item Eine Architektur-Governance entwerfen, die leicht, aber wirksam ist \textendash{} und Entscheidungen unter Unsicherheit dokumentieren.
\end{itemize}
\end{infobox}

\vspace{1cm}
\noindent{\color{lpAccent}\rule{\linewidth}{0.6pt}}\\[4pt]
{\footnotesize\sffamily\color{lpGray}Lernplatt \textperiodcentered{} by Can Siebert \textperiodcentered{} Kurs 71 (Dig 42) \textperiodcentered{} Tag 10 \textperiodcentered{} Aufgabenheft \textperiodcentered{} \textcopyright{} 2026}

\end{document}
